%==============================================================================
% فصل سوم: مواد و روش کار
% جزئیات باید به‌اندازه‌ای باشد که مطالعه قابل تکرار باشد.
%==============================================================================
\ChapterStart
\chapter{مواد و روش کار}
\label{ch:methods}

\section{نوع پژوهش}
\label{sec:study-design}
پژوهش حاضر یک مطالعه مداخله‌ای از نوع کارآزمایی کنترل شده تصادفی‌ با دو بازوی موازی بود که با هدف مقایسه اثربخشی آموزش خودمراقبتی به روش کارگاهی و آموزش کارگاهی تلفیق‌شده با خرده‌یادگیری بر عملکرد خودمراقبتی افراد مبتلا به پرفشاری خون انجام شد. در این مطالعه، هر دو گروه آموزش حضوری کارگاهی را دریافت کردند و مداخله خرده‌یادگیری فقط به برنامه آموزشی یکی از گروه‌ها افزوده شد. بنابراین، گروه دریافت‌کننده آموزش کارگاهی به‌عنوان گروه مقایسه فعال و گروه دریافت‌کننده آموزش کارگاهی همراه با خرده‌یادگیری به‌عنوان گروه مداخله در نظر گرفته شد. پیامدهای مطالعه شامل عملکرد خودمراقبتی و میانگین فشار خون خانگی بود که در مرحله پایه و یک ماه پس از پایان مداخله اندازه‌گیری شد.

\section{جامعه، محیط و نمونه پژوهش}
\label{sec:population}
جامعه پژوهش را بیماران ۳۰ تا ۷۰ ساله مبتلا به پرفشاری خون تشکیل دادند که به درمانگاه قلب بیمارستان شهید مدنی خرم‌آباد مراجعه کرده بودند. محیط جذب نمونه، درمانگاه قلب بیمارستان شهید مدنی و محیط اجرای جلسات حضوری، کلاس‌های آموزشی بیمارستان مذکور بود. نمونه پژوهش از میان مراجعانی انتخاب شد که تشخیص قطعی پرفشاری خون آنان به‌وسیله متخصص قلب تایید شده بود و سایر معیارهای ورود را داشتند.
\par
در آغاز فرایند جذب، مراجعان از نظر معیارهای ورود و خروج غربال شدند. سپس اهداف مطالعه، مراحل اجرا، نوع مداخله آموزشی، شیوه سنجش پیامدها و اختیاری بودن مشارکت برای افراد واجد شرایط توضیح داده شد. افرادی که پس از دریافت توضیحات لازم، رضایت‌نامه آگاهانه کتبی را امضا کردند، وارد مطالعه شدند. برای هر شرکت‌کننده یک کد شناسایی اختصاص یافت تا اطلاعات پژوهشی بدون درج نام و مشخصات هویتی ثبت و نگهداری شود.

\section{روش نمونه‌گیری}
\label{sec:sampling}
بیماران مراجعه‌کننده به درمانگاه قلب بیمارستان شهید مدنی خرم‌آباد از نظر شرایط ورود به مطالعه بررسی شدند و افراد واجد شرایط پس از اعلام رضایت وارد پژوهش شدند. پس از تکمیل اندازه‌گیری‌های پایه، شرکت‌کنندگان با استفاده از تصادفی‌سازی بلوکی و بلوک‌هایی با اندازه متغیر به یکی از دو گروه آموزش کارگاهی یا آموزش کارگاهی تلفیق‌شده با خرده‌یادگیری تخصیص یافتند. استفاده از اندازه‌های متغیر بلوک با هدف حفظ تعادل تعداد افراد در دو گروه در طول جذب و کاهش امکان پیش‌بینی توالی تخصیص انجام شد. تولید و اجرای توالی تخصیص به فردی خارج از تیم پژوهش سپرده شد.

\section{حجم نمونه و روش محاسبه آن}
\label{sec:sample-size}
حجم نمونه با استفاده از فرمول مقایسه میانگین دو جامعه مستقل محاسبه شد. فرمول مورد استفاده به صورت زیر بود:

\[
n=\frac{\left(Z_{1-\frac{\alpha}{2}}+Z_{1-\beta}\right)^2\left(\sigma_1^2+\sigma_2^2\right)}{\left(\mu_1-\mu_2\right)^2}
\]

در این فرمول، \(n\) تعداد نمونه مورد نیاز در هر گروه، \(\alpha\) احتمال خطای نوع اول، \(1-\beta\) توان آزمون، \(\sigma_1^2\) و \(\sigma_2^2\) واریانس پیامد در دو گروه و \(\mu_1-\mu_2\) تفاوت مورد انتظار میانگین پیامد بین دو گروه بود. با در نظر گرفتن خطای نوع اول برابر با $0.05$، توان ۸۰ درصد و احتمال ریزش ۱۰ درصد، حداقل حجم نمونه ۵۵ نفر برای هر گروه و در مجموع ۱۱۰ نفر تعیین شد.

\input{tables/table-4-1}

\section{معیارهای ورود و خروج}
\label{sec:eligibility}
معیارهای ورود شامل سن ۳۰ تا ۷۰ سال، تشخیص جدید و قطعی پرفشاری خون به‌وسیله متخصص قلب، توانایی خواندن و نوشتن، امکان حضور در جلسات آموزشی، دسترسی به تلفن همراه هوشمند و اینترنت، توانایی استفاده از محتوای آموزشی ارسالی و توانایی کار با دستگاه فشارسنج خانگی بود. برای کاهش خطای سنجش، توانایی عملی شرکت‌کنندگان در اندازه‌گیری فشار خون پیش از شروع گردآوری داده‌های پایه بررسی شد.
\par
ابتلا به سایر بیماری‌های مزمن، غیبت در بیش از یک جلسه آموزشی و نداشتن تمایل به ادامه همکاری از معیارهای خروج از مطالعه بود. خروج داوطلبانه شرکت‌کننده در هر مرحله امکان‌پذیر بود و انصراف از پژوهش هیچ اثری بر دریافت مراقبت و درمان معمول وی نداشت.

\section{ابزار گردآوری اطلاعات}
\label{sec:instruments}
برای گردآوری داده‌ها از فرم مشخصات جمعیت‌شناختی و بالینی، پرسشنامه خودمراقبتی پرفشاری خون و فرم ثبت فشار خون خانگی استفاده شد. فرم مشخصات جمعیت‌شناختی و بالینی برای ثبت متغیرهایی مانند سن، جنسیت، سطح تحصیلات، شغل و اطلاعات بالینی مرتبط به‌کار رفت.
\par
عملکرد خودمراقبتی با نسخه فارسی پرسشنامه خودمراقبتی پرفشاری خون  «\lr{HTN-SCP}» سنجیده شد. این پرسشنامه یک ابزار اندازه‌گیری استاندارد است که بر اساس مدل رایگل خودمراقبتی در بیماری‌های مزمن طراحی شده و سه بُعد عملکرد، رفتار و شناخت را می‌سنجد \cite{han2014hscp}.
\par
بُعد عملکرد شامل ۲۰ گویه است و رفتارهای واقعی بیمار در حوزه‌های پایش، دارو، تغذیه، فعالیت بدنی و مدیریت استرس را می‌سنجد. بُعد رفتار شامل ۲۰ گویه است و باورها، نگرش‌ها و بازتاب‌های بیماران درباره انجام رفتارهای خودمراقبتی را ارزیابی می‌کند. بُعد شناخت شامل ۲۰ گویه است و دانش، آگاهی و درک بیمار از پرفشاری خون و رفتارهای خودمراقبتی را می‌سنجد. مجموع پرسشنامه شامل ۶۰ گویه است.
\par
پرسشنامه بر اساس طیف لیکرت پنج‌درجه‌ای نمره‌گذاری می‌شود، به طوری که هر گویه نمره ۱ تا ۵ می‌گیرد و نمره کل بین ۶۰ تا ۳۰۰ قرار می‌گیرد. نمره بالاتر نشان‌دهنده عملکرد بهتر خودمراقبتی است.
روایی و پایایی نسخه فارسی این پرسشنامه توسط قانعی قشلاقی و همکاران بررسی و تأیید شده است. ضریب آلفای کرونباخ کل ابزار $0.86$ گزارش شده است، که نشان‌دهنده پایایی قابل‌قبول است. روایی محتوایی و روایی سازه نیز توسط این مطالعه تأیید شده است. این پرسشنامه ابزار انتخابی پژوهش حاضر است، چرا که به‌طور خاص برای بیماران مبتلا به پرفشاری خون طراحی شده و سه بُعد عملکرد، رفتار و شناخت را هم‌زمان می‌سنجد. استفاده از این پرسشنامه امکان سنجشی دقیق از ابعاد مختلف خودمراقبتی را فراهم می‌آورد و مقایسه پیش و پس از مداخله را میسر می‌سازد \cite{ghaneigheshlagh2019persian}.
\par
چند ویژگی این پرسشنامه آن را ابزاری مناسب برای پژوهش حاضر می‌سازد. نخست، تخصصی بودن برای پرفشاری خون؛ برخلاف پرسشنامه‌های عمومی خودمراقبتی، این پرسشنامه بر بازتاب‌های رفتاری مشخص در پرفشاری خون تمرکز دارد. دوم، سه‌گانه بودن و پوشش بُعد شناختی؛ این ویژگی امکان تحلیل جداگانه پیشاوردهای شناختی و رفتاری را فراهم می‌سازد، که با یافته‌هایی که خرده‌یادگیری بر پیشاوردهای شناختی اثربخشی قابل‌اعتمادتری دارد و بر رفتار گسست دارد، هم‌خوان است. سوم، اعتبارسنجی شده در فارسی؛ این امر امکان کاربست را در جمعیت ایرانی بدون نیاز به کار روان‌سنجی بیشتر فراهم می‌آورد.
نمره حاصل از پرسشنامه به‌عنوان یکی از پیامدهای اصلی مطالعه در مرحله پیش‌آزمون و پس آزمون (یک‌ماه پس از مداخله) ثبت شد.
\par
برای اندازه‌گیری فشار خون، از دستگاه فشارسنج دیجیتال استاندارد خانگی استفاده شد. پیش از شروع ثبت خانگی، شیوه صحیح استفاده از دستگاه به‌صورت عملی آموزش داده شد و هر شرکت‌کننده یک بار تمام مراحل اندازه‌گیری را در حضور پژوهشگر انجام داد. پژوهشگر با استفاده از سیاهه عملکرد، درستی اجرای مراحل را بررسی کرد. برای تعیین مقدار پایه، شرکت‌کنندگان فشار خون خود را به مدت یک هفته، روزانه در چهار نوبت شامل دو بار در صبح و دو بار در شب اندازه‌گیری و در فرم مربوط ثبت کردند. میانگین اندازه‌گیری‌های ثبت‌شده به‌عنوان فشار خون پایه در نظر گرفته شد. همین روش ثبت خانگی در پیگیری‌های پس از مداخله نیز به‌کار رفت.

\section{روش گردآوری اطلاعات و اجرای مداخله}
\label{sec:procedure}
پس از غربالگری و اخذ رضایت آگاهانه، اطلاعات جمعیت‌شناختی و بالینی ثبت شد. سپس آموزش عملی اندازه‌گیری فشار خون خانگی انجام گرفت و شرکت‌کنندگان اندازه‌گیری‌های پایه را طی یک هفته ثبت کردند. پیش از شروع آموزش نیز پرسشنامه خودمراقبتی پرفشاری خون تکمیل شد. پس از پایان سنجش پایه، تخصیص تصادفی انجام گرفت و برنامه آموزشی چهار هفته‌ای آغاز شد.
\par
آموزش بر اساس الگوی طراحی آموزشی \lr{ADDIE} و در پنج مرحله تحلیل، طراحی، توسعه، اجرا و ارزشیابی انجام شد. در مرحله تحلیل، ویژگی‌های فردی و اجتماعی شرکت‌کنندگان و وضعیت اولیه آگاهی و عملکرد خودمراقبتی آنان بررسی و شکاف‌های دانشی، نگرشی و مهارتی در زمینه شناخت بیماری، مصرف دارو، تغذیه، فعالیت بدنی، مدیریت تنش و پایش فشار خون شناسایی شد. در مرحله طراحی، اهداف رفتاری قابل اندازه‌گیری، ساختار جلسات حضوری، روش‌ تدریس و ساختار محتوای خرده‌یادگیری تعیین شد. در مرحله توسعه، اسلایدهای آموزشی، مطالب راهنما و هشت ویدیوی کوتاه خرده‌یادگیری آماده شد. در مرحله اجرا، جلسات مربوط به آموزش کارگاهی اجرا و ارسال برنامه‌ریزی‌شده میکرو ویدیوها انجام گرفت. در مرحله ارزشیابی، حضور شرکت‌کنندگان، مشاهده ویدیوها، بازخوردهای کوتاه آنان و تغییر پیامدهای مطالعه پایش شد.
\par
هر دو گروه در 4 جلسه حضوری هفتگی شرکت کردند. هر جلسه ۶۰ دقیقه طول کشید و جلسات در چهار جمعه متوالی در کلاس‌های آموزشی بیمارستان شهید مدنی خرم‌آباد برگزار شد. محتوای آموزشی با استفاده از سخنرانی کوتاه، پرسش و پاسخ، بحث گروهی و تمرین عملی ارائه شد. در گروه آموزش کارگاهی، آموزش مذکور به صورت کارگاهی و حضوری ارائه شد اما، گروه آموزش کارگاهی تلفیق‌شده با خرده‌یادگیری، علاوه بر جلسات حضوری، 8 ویدیوی آموزشی سه‌دقیقه‌ای دریافت کردند. ویدیوها دو بار در هفته از طریق پیام‌رسان تلفن همراه ارسال و دریافت تایید مشاهده هر ویدیو برای پایش پایبندی آموزشی در نظر گرفته شد. ساختار مداخله در جدول ۳ـ۱ ارائه شده است.

\begin{table}[htbp]
\centering
\caption{ساختار و محتوای مداخله آموزشی در دو گروه پژوهش}
\label{tab:intervention}
{\small
\setlength{\tabcolsep}{3pt}
\renewcommand{\arraystretch}{1.2}
\begin{tabularx}{0.985\textwidth}{p{0.10\textwidth} X p{0.18\textwidth} p{0.22\textwidth}}
\hline
\textbf{مرحله} & \textbf{محتوای آموزشی} & \textbf{گروه آموزش کارگاهی} & \textbf{گروه کارگاهی تلفیق‌شده با خرده‌یادگیری} \\
\hline
جلسه اول & آشنایی با پرفشاری خون و عوارض آن، مفهوم خودمراقبتی و روش صحیح اندازه‌گیری فشار خون در منزل & به صورت حضوری (۶۰ دقیقه‌) & به صورت حضوری (۶۰ دقیقه‌) همراه با میکروویدیوهای ۱ و ۲ \\
جلسه دوم & تغذیه سالم، کاهش نمک و چربی، آشنایی با الگوی غذایی «\lr{DASH}» و کنترل وزن & به صورت حضوری (۶۰ دقیقه‌) & به صورت حضوری (۶۰ دقیقه‌) همراه با میکروویدیوهای 3 و 4 \\
جلسه سوم & فعالیت بدنی منظم، ورزش متناسب با توان فردی، مدیریت تنش و راهبردهای آرام‌سازی & به صورت حضوری (۶۰ دقیقه‌) &به صورت حضوری (۶۰ دقیقه‌) همراه با میکروویدیوهای 5 و 6 \\
جلسه چهارم & داروهای رایج پرفشاری خون، پایبندی دارویی، مدیریت خودمراقبتی در زندگی روزمره و جمع‌بندی & به صورت حضوری (۶۰ دقیقه‌) & به صورت حضوری (۶۰ دقیقه‌) همراه با میکروویدیوهای 7 و 8 \\
\hline
\end{tabularx}
}
\end{table}


\subsection{میکروویدیوهای آموزشی (خرده یادگیری)}

میکروویدیوی نخست به معرفی کوتاه پرفشاری خون و اهمیت کنترل آن اختصاص داشت و ویدیوی دوم روش صحیح اندازه‌گیری فشار خون در منزل را به‌صورت خلاصه و تصویری نشان می‌داد. میکروویدیوهای سوم و چهارم به‌ترتیب اصول کاهش نمک و چربی و نکات عملی انتخاب و آماده‌سازی غذای مناسب را پوشش می‌دادند. میکروویدیوهای پنجم و ششم بر شروع فعالیت بدنی متناسب با توان فردی و مرور فواید قلبی‌عروقی فعالیت منظم تمرکز داشتند. میکروویدیوی هفتم اهمیت مصرف منظم داروهای تجویزشده و پرهیز از قطع خودسرانه درمان را یادآوری می‌کرد و نهایتا میکروویدیوی هشتم جمع‌بندی نکات خودمراقبتی و تشویق به ادامه ثبت فشار خون و حفظ سبک زندگی سالم را دربر می‌گرفت.
\par
ارزشیابی فرایندی مداخله با ثبت حضور و غیاب جلسات، پایش مشاهده میکروویدیوها و دریافت بازخوردهای کوتاه شرکت‌کنندگان انجام شد. ارزشیابی پیامدها در دو نوبت شامل پیش‌آزمون و یک ماه پس از پایان مداخله صورت گرفت. در هر نوبت، عملکرد خودمراقبتی و میانگین فشار خون خانگی بیماران ثبت شد. زمان‌بندی سنجش متغیرها در جدول ۳ـ۲ ارائه شده است.

\begin{table}[htbp]
\centering
\caption{زمان‌بندی گردآوری داده‌ها در مراحل پژوهش}
\label{tab:measurements}
{\small
\setlength{\tabcolsep}{3pt}
\renewcommand{\arraystretch}{1.2}
\begin{tabularx}{0.985\textwidth}{X *{2}{>{\centering\arraybackslash}p{0.22\textwidth}}}
\hline
\textbf{متغیر یا فعالیت} & \textbf{پیش از مداخله} & \textbf{یک ماه پس از مداخله} \\
\hline
ثبت مشخصات جمعیت‌شناختی و بالینی & $\checkmark$ & ــ \\
تکمیل پرسشنامه خودمراقبتی پرفشاری خون & $\checkmark$ & $\checkmark$ \\
ثبت فشار خون خانگی طی یک هفته & $\checkmark$ & $\checkmark$ \\
بررسی حضور در جلسات کارگاهی & ــ & $\checkmark$ \\
پایش مشاهده میکروویدیوها & ــ & $\checkmark$ \\
\hline
\end{tabularx}
}
\end{table}



جلسات در روزهای جمعه برنامه‌ریزی شد تا حداقل تداخل با مشغله روزمره بیماران داشته باشد. در زمان ثبت‌نام نیز آموزش اولیه استفاده از فشارسنج دیجیتال در اختیار گروه های مداخله قرار گرفت. آموزش عملی اندازه‌گیری فشار خون و بررسی عملکرد شرکت‌کنندگان پیش از ثبت خانگی برای کاهش خطای اندازه‌گیری انجام شد.

\section{روش‌های آماری و تجزیه‌وتحلیل داده‌ها}
\label{sec:statistics}
داده‌ها پس از بررسی اولیه، با نرم‌افزار آماری \lr{SPSS}؛ نسخه ۲۶ تحلیل شدند. برای توصیف متغیرهای کیفی از فراوانی و درصد و برای توصیف متغیرهای کمی از میانگین و انحراف معیار استفاده شد. سطح معناداری آزمون‌های آماری کمتر از ۰/۰۵ در نظر گرفته شد.
\par

ارزیابی همگنی دو گروه در مرحله پایه و برای متغیرهای کمی با کمک آزمون تی مستقل و برای متغیرهای کیفی با استفاده  آزمون کای‌دو انجام شد.. تغییرات درون گروهی با آزمون تی زوجی و مقایسه­های  بین‌گروهی پیامدها و برآورد اثر افزوده خرده‌یادگیری، با کمک آزمون تحلیل کوواریانس انجام گرفت. به این ترتیب که مقدار پیش‌آزمون هر پیامد به‌عنوان متغیر کمکی در مدل وارد و اثر گروه پس از کنترل مقدار پایه سنجیده شد.

\section{ملاحظات اخلاقی}
\label{sec:ethics}
پژوهش پس از دریافت تاییدیه کمیته اخلاق در پژوهش دانشگاه علوم پزشکی تهران اجرا شد. پیش از ورود افراد به مطالعه، هدف، روش اجرا، منافع مورد انتظار، تعهدات شرکت‌کننده و حق خروج از مطالعه برای آنان توضیح داده شد و رضایت‌نامه آگاهانه کتبی دریافت شد. مشارکت کاملا داوطلبانه بود و شرکت‌کنندگان می‌توانستند در هر مرحله بدون تاثیر بر روند درمان معمول خود از ادامه همکاری انصراف دهند.
\par
برای رعایت محرمانگی، داده‌ها با کدهای شناسایی فاقد نام ثبت شدند و اطلاعات هویتی در گزارش پژوهش انعکاس نیافت. دسترسی به داده‌ها به افراد مجاز پژوهش محدود شد. مداخله مطالعه ماهیت آموزشی داشت و بر اصلاح سبک زندگی، پایش خانگی فشار خون و تقویت پایبندی به درمان تجویزشده متمرکز بود. در تمام آموزش‌ها تاکید شد که شرکت‌کنندگان از تغییر یا قطع خودسرانه داروها خودداری کنند و درمان خود را مطابق نظر پزشک ادامه دهند.
